\section{Governance Implications, Conditional on the Bounds}

Every implication below is conditional on the bounds of \S4; none survives
transposition into an unconditional policy claim.

\begin{enumerate}\itemsep1pt
\item \textbf{Targeted dismantling is the only lever that moves the structure.}
Random removal needs about 60\% of the financial graph to halve its giant
component whereas degree-targeted removal needs 2.5\% (E6-fin). \emph{Condition:}
the giant component is a weak proxy at 47/49 saturated steps, so this is about
relative leverage, not absolute effect size.
\item \textbf{The illicit financial subgraph is a periphery, so act accordingly.}
Removing 100\% of illicit-labelled addresses leaves the component at 98.88\%
(E6-fin). \emph{Condition:} address-level enforcement is disruption of service
rather than damage to the network, and expecting network-level effects from it
is a category error our data bounds.
\item \textbf{The fast layer is where resilience lives.} 13.88\% of darkened
hosts return within a day at 6.2$\times$ chance (E7-inf) and median
infrastructure lifetime is 1.36 days (E0). \emph{Condition:} takedown cadence,
not takedown existence, is the operative variable at this layer.
\item \textbf{Domain infrastructure is the durable substrate.} Domain-hosted
infrastructure returns 1.31$\times$ as often as bare IP (80.77\% vs 61.49\%,
p \textless{} 0.001). \emph{Condition:} interventions that raise the cost of
domain re-registration attack the slower half of the churn cycle, and the feed's
scanner admixture is not separable.
\item \textbf{Displacement is directional: single-market seizures relocate rather
than remove.} Exposed vendors surface on Apollon at +10.5 pp above controls and
on Empire at +0.0, a zero-denominator cell that is unmeasurable rather than a
measured zero (E1-v3). \emph{Condition:} the corrected figures are two
non-significant point estimates, so the directional design and not the magnitude
is the finding, and +10.5 pp is an upper bound on substitution.
\end{enumerate}

\section{Limitations}

We consolidate, rather than disperse, the weaknesses a reviewer would find.

(i) \textbf{Causal magnitudes are not identified.} E3's design failed at matching
(overlap = 0.000), and with one ecosystem-wide shock no placebo or matching
machinery recovers an effect size; only H4 survives, and H4's support is the
confirmation of an absence, which is weaker evidence than a positive finding.
This is not a null result about the mechanism; it is the demonstration that this
data cannot support effect-size estimation in any design we could build.

(ii) \textbf{Historical windows.} The dated parts of the social and financial
layers cover roughly 2014--2017: the AlphaBay listings span 2014-12-21 to
2017-01-28, Agora 2014--2015, and the financial layer's ordinal axis maps to
approximately 2016-01-01 to 2017-11-03. Two social-layer panels (Kilos,
drug-listing) carry no date column at all, so no window can be stated for them,
and the current layer is infrastructure only.

(iii) \textbf{Censoring.} The social-layer outcome is first recorded feedback
rather than first transaction, and no seller key exists, so the event rate is a
lower bound and the standard errors are optimistic (E4). Both defects push every
duration estimate in the same direction, which is why the two medians are always
reported as a pair.

(vi) \textbf{Label provenance.} Elliptic++ labels are dataset-internal and
unverifiable against external ground truth, and the E5-fin gradient is equally
consistent with collection economics; the metrics that rest on this labelling
layer inherit this base.

Dropped limitations (metric-validity detail, external validation, novelty
exhaustiveness) are retained in Supplementary Notes S2 and S5.
